\section{Результаты экспериментов}\label{sec:results}

\subsection{Репликация базового эксперимента}

На рис.~\ref{fig:base} представлена карта обучаемости для базовой конфигурации
(tanh, $|\mathcal{D}|=1$, полный батч, разрешение $512\times512$). Слева ---
карта четырёх режимов, справа --- непрерывная сине-красная карта в стиле
предыдущей НИР, где насыщенность синего пропорциональна сумме нормализованных
loss по итерациям, а красного --- обратной величине.

\begin{figure}[htbp]
  \centering
  \includegraphics[width=\linewidth]{f01_base_tanh.png}
  \caption{Базовая карта обучаемости в плоскости $(\eta_0,\eta_1)$, активация
  tanh, разрешение $512\times512$, $T=400$ итераций.}
  \label{fig:base}
\end{figure}

Качественная картина воспроизводит результат предыдущей НИР: существуют
протяжённые связные области успешного обучения, граница между ними и областью
расходимости не является гладкой кривой, а оптимальные по скорости
конфигурации располагаются вблизи границы. Существенное отличие от предыдущей
работы состоит в том, что значительную часть карты занимает режим ограниченных
незатухающих колебаний (доля \DBaseTanhConv{} приходится на сходимость), который
при бинарном критерии предыдущей НИР был бы отнесён к сходимости.

Измеренная размерность границы составила
$D_{\mathrm{box}} = \DBaseTanh \pm \DBaseTanhS$.

Сравнение активаций осложняется тем, что порог устойчивости у них различен
(см.~табл.~\ref{tab:act}). В общем окне $\eta\in[10^{-2},10]$ для ReLU при
$|\mathcal{D}|=1$ доля сходимости равна $0.99$, а для линейной сети, наоборот,
$0.07$: граница выходит за пределы окна. Поэтому дополнительно выполнены
прогоны в окнах, центрированных на собственном пороге каждой конфигурации, и
прогоны на общем наборе из 32 примеров, где сравнение активаций корректно:
tanh даёт $\DTanhMb$, ReLU --- $\DReluMbLow$.

\begin{table}[htbp]
  \centering
  \caption{Размерность границы для разных активаций и окон сканирования.
  Столбцы «сход.», «осц.», «расх.», «NaN» --- доли соответствующих режимов.}
  \label{tab:act}
  \small
  \begin{tabular}{llccccc}
\toprule
Активация & Окно по $\eta$ & $D_{\mathrm{box}}$ & сход. & осц. & расх. & NaN\\
\midrule
tanh & $[10^{-2},10]$, $|\mathcal{D}|=1$, $512^2$ & $1.688 \pm 0.035$ & 0.49 & 0.19 & 0.06 & 0.26\\
ReLU & $[10^{-2},10]$, $|\mathcal{D}|=1$ & $1.160 \pm 0.008$ & 0.99 & 0.01 & 0.00 & 0.00\\
линейная & $[10^{-2},10]$, $|\mathcal{D}|=1$ & $1.512 \pm 0.067$ & 0.07 & 0.02 & 0.00 & 0.91\\
линейная & $[10^{-4},10^{-1}]$, $|\mathcal{D}|=1$ & $1.565 \pm 0.033$ & 0.75 & 0.15 & 0.00 & 0.10\\
tanh & $[10^{-1},10^{3}]$, $|\mathcal{D}|=32$ & $1.433 \pm 0.050$ & 0.38 & 0.07 & 0.02 & 0.54\\
ReLU & $[10^{-3},10]$, $|\mathcal{D}|=32$ & $1.282 \pm 0.045$ & 0.47 & 0.11 & 0.00 & 0.43\\
ReLU & $[10^{-1},10^{3}]$, $|\mathcal{D}|=32$ & $1.074 \pm 0.010$ & 0.06 & 0.00 & 0.00 & 0.94\\
\bottomrule
\end{tabular}

\end{table}

Рис.~\ref{fig:act} показывает карты для трёх активаций в общем окне; хорошо
видно, что для ReLU граница в это окно не попадает.

\begin{figure}[htbp]
  \centering
  \includegraphics[width=\linewidth]{f02_activations.png}
  \caption{Карты режимов для активаций tanh, ReLU и линейной в общем окне
  $\eta\in[10^{-2},10]$.}
  \label{fig:act}
\end{figure}

Линейность зависимости $\log N(\varepsilon)$ от $\log\varepsilon$,
подтверждающая применимость оценки~\eqref{eq:box}, показана на
рис.~\ref{fig:boxcount}.

\begin{figure}[htbp]
  \centering
  \includegraphics[width=0.62\linewidth]{f03_boxcounting.png}
  \caption{Box-counting: зависимость числа непустых ячеек от размера ячейки
  в логарифмических координатах для базовой карты и для сечения
  $(\eta,\beta)$.}
  \label{fig:boxcount}
\end{figure}

\subsection{Сечение $(\eta,\beta)$: инерция}

Центральный результат работы представлен на рис.~\ref{fig:etabeta}: сечение
пространства гиперпараметров по шагу обучения и коэффициенту инерции при
общем для обоих слоёв $\eta$.

\begin{figure}[htbp]
  \centering
  \includegraphics[width=\linewidth]{f04_eta_beta.png}
  \caption{Сечение $(\eta,\beta)$ для SGD с инерцией, разрешение
  $512\times512$. Справа --- извлечённая скелетонизированная граница
  множества сходимости.}
  \label{fig:etabeta}
\end{figure}

Измеренная размерность $D_{\mathrm{box}} = \DEtaBeta \pm \DEtaBetaS$ совпадает
в пределах погрешности с размерностью базовой карты $(\eta_0,\eta_1)$,
равной $\DBaseTanh \pm \DBaseTanhS$. Таким образом, добавление инерции не
делает границу обучаемости гладкой: первая из проверяемых гипотез
подтверждается.

Структура карты согласуется с теорией раздела~\ref{sec:theory}. При малых
$\eta$ обучение сходится при любом $\beta$; при росте $\eta$ система сначала
входит в область ограниченных колебаний, и лишь затем --- в область
переполнения. Граница области сходимости смещается влево с ростом $\beta$,
что отражает увеличение эффективного шага: в стационарном режиме накопленная
скорость в схеме~\eqref{eq:hb} превышает мгновенный градиент в $1/(1-\beta)$
раз. При $\beta\gtrsim0.8$ область незатухающих колебаний практически исчезает,
уступая место резкому переходу от сходимости к переполнению.

\subsection{Затухание весов и размер батча}

Результаты для сечения $(\eta,\lambda_{wd})$ и для мини-батчей сведены в
табл.~\ref{tab:sections} и~\ref{tab:mb}, карты приведены на
рис.~\ref{fig:wd}--\ref{fig:mb}.

\begin{table}[htbp]
  \centering
  \caption{Размерность границы в основных сечениях (разрешение $256\times256$,
  кроме отмеченных отдельно).}
  \label{tab:sections}
  \small
  \input{tables/tab_sections.tex}
\end{table}

\begin{figure}[htbp]
  \centering
  \begin{subfigure}{0.46\linewidth}
    \includegraphics[width=\linewidth]{f07_eta_wd.png}
    \caption{сечение $(\eta,\lambda_{wd})$}
  \end{subfigure}\hfill
  \begin{subfigure}{0.52\linewidth}
    \includegraphics[width=\linewidth]{f08_eta_batch.png}
    \caption{сечение $(\eta,B)$}
  \end{subfigure}
  \caption{Влияние затухания весов и размера батча.}
  \label{fig:wd}
\end{figure}

Затухание весов заметно сглаживает границу: $D_{\mathrm{box}} = \DEtaWd \pm \DEtaWdS$
против $\DSgd \pm \DSgdS$ для карты $(\eta_0,\eta_1)$ при том же разрешении.
Механизм прозрачен: слагаемое $\lambda_{wd}W$ добавляет к гессиану
изотропный сдвиг $\lambda_{wd}I$, что при больших $\lambda_{wd}$ подавляет
различие между направлениями разной кривизны и делает порог устойчивости
почти не зависящим от точки в пространстве параметров.

Для сечения $(\eta,B)$ карта режимов построена (рис.~\ref{fig:wd}б), однако
размерность по нему не оценивалась: полоса $512\times32$ не даёт достаточного
числа масштабов $\varepsilon$ для устойчивой регрессии~\eqref{eq:box}, и
результат оказывается неопределённым. Вместо этого влияние размера батча
измерено по полноценным картам $(\eta_0,\eta_1)$ при фиксированных $B$.

\begin{table}[htbp]
  \centering
  \caption{Влияние размера мини-батча ($|\mathcal{D}|=32$, tanh).}
  \label{tab:mb}
  \small
  \begin{tabular}{lccccc}
\toprule
Конфигурация & $D_{\mathrm{box}}$ & сход. & осц. & расх. & NaN\\
\midrule
$B=1$ & $1.722 \pm 0.074$ & 0.17 & 0.09 & 0.13 & 0.61\\
$B=4$ & $1.568 \pm 0.050$ & 0.30 & 0.08 & 0.09 & 0.53\\
$B=32$ (полный батч) & $1.433 \pm 0.050$ & 0.38 & 0.07 & 0.02 & 0.54\\
\bottomrule
\end{tabular}

\end{table}

\begin{figure}[htbp]
  \centering
  \includegraphics[width=\linewidth]{f09_minibatch.png}
  \caption{Карты обучаемости при разных размерах мини-батча.}
  \label{fig:mb}
\end{figure}

Наблюдается монотонный тренд: $B=1$ даёт $\DMbOne$, $B=4$ --- $\DMbFour$,
полный батч $B=32$ --- $\DMbFull$. Уменьшение батча увеличивает изрезанность
границы. Это согласуется с интерпретацией~\cite{kong2020}: стохастичность
выборки играет роль дополнительного мелкомасштабного возмущения целевой
функции, а по результатам~\cite{liu2024} именно мелкомасштабная невыпуклость
и порождает фрактальную границу. Вторая гипотеза, таким образом,
подтверждается для размера батча, но в части затухания весов --- лишь при
достаточно больших $\lambda_{wd}$ (см.~ниже раздел~\ref{sec:3d}).

\subsection{Сравнение оптимизаторов}

Результаты сравнения четырёх методов приведены в табл.~\ref{tab:opt} и на
рис.~\ref{fig:opt}.

\begin{table}[htbp]
  \centering
  \caption{Сравнение оптимизаторов и сечений Adam (разрешение $256\times256$).}
  \label{tab:opt}
  \small
  \begin{tabular}{lccccc}
\toprule
Конфигурация & $D_{\mathrm{box}}$ & сход. & осц. & расх. & NaN\\
\midrule
SGD & $1.633 \pm 0.041$ & 0.49 & 0.19 & 0.06 & 0.26\\
SGD+momentum ($\beta=0.9$) & $1.037 \pm 0.011$ & 0.82 & 0.00 & 0.02 & 0.16\\
Adam & $1.322 \pm 0.054$ & 0.82 & 0.05 & 0.14 & 0.00\\
RMSprop & $1.508 \pm 0.068$ & 0.20 & 0.54 & 0.26 & 0.00\\
Adam, сечение $(\eta,\beta_1)$ & $1.698 \pm 0.071$ & 0.48 & 0.52 & 0.00 & 0.00\\
Adam, сечение $(\eta,\beta_2)$ & $1.727 \pm 0.073$ & 0.28 & 0.72 & 0.00 & 0.00\\
\bottomrule
\end{tabular}

\end{table}

\begin{figure}[htbp]
  \centering
  \includegraphics[width=\linewidth]{f10_optimizers.png}
  \caption{Карты обучаемости в плоскости $(\eta_0,\eta_1)$ для четырёх
  оптимизаторов. Диапазоны по осям различны и подобраны под собственный порог
  каждого метода.}
  \label{fig:opt}
\end{figure}

В плоскости $(\eta_0,\eta_1)$ Adam действительно даёт более гладкую границу,
чем SGD: $\DAdam$ против $\DSgd$. Однако наименьшую размерность даёт вовсе не
адаптивный метод, а SGD с фиксированной инерцией $\beta=0.9$: $\DMomentum$.
RMSprop занимает промежуточное положение ($\DRmsprop$).

Более того, если сечь пространство гиперпараметров вдоль осей самого Adam,
картина меняется на противоположную (рис.~\ref{fig:adamsec}): сечение
$(\eta,\beta_1)$ даёт $\DAdamBOne$, а $(\eta,\beta_2)$ --- $\DAdamBTwo$, то
есть выше, чем у SGD в базовом сечении. Гипотеза о сглаживающем действии
адаптивных методов, таким образом, подтверждается лишь частично и требует
переформулировки, к чему мы вернёмся в разделе~\ref{sec:discussion}.

\begin{figure}[htbp]
  \centering
  \includegraphics[width=\linewidth]{f11_adam_sections.png}
  \caption{Сечения пространства гиперпараметров Adam по $\beta_1$ и $\beta_2$.}
  \label{fig:adamsec}
\end{figure}

Отдельного внимания заслуживает структура карт Adam и RMSprop: на них хорошо
видны горизонтальные полосы, почти не зависящие от $\eta_0$. Это прямое
следствие покоординатной нормализации: величина шага по каждому весу равна
приблизительно $\eta$ независимо от масштаба градиента, поэтому динамика
разделяется по слоям, и карта становится почти сепарабельной. Именно эта
сепарабельность, а не подавление хаоса как такового, и снижает измеряемую
размерность границы в плоскости $(\eta_0,\eta_1)$.

\subsection{Случайные и структурированные данные}

\begin{table}[htbp]
  \centering
  \caption{Влияние структуры данных ($|\mathcal{D}|=16$, полный батч, tanh).}
  \label{tab:data}
  \small
  \begin{tabular}{lccccc}
\toprule
Конфигурация & $D_{\mathrm{box}}$ & сход. & осц. & расх. & NaN\\
\midrule
случайные данные, $|D|=16$ & $1.284 \pm 0.070$ & 0.39 & 0.02 & 0.02 & 0.57\\
digits (PCA-16), $|D|=16$ & $1.494 \pm 0.073$ & 0.34 & 0.04 & 0.02 & 0.60\\
\bottomrule
\end{tabular}

\end{table}

\begin{figure}[htbp]
  \centering
  \includegraphics[width=\linewidth]{f12_data.png}
  \caption{Границы обучаемости на случайных данных и на наборе digits.}
  \label{fig:data}
\end{figure}

При переходе от случайных данных к структурированным размерность возрастает:
$\DSynth$ против $\DDigits$ (табл.~\ref{tab:data}, рис.~\ref{fig:data}).
Правдоподобное объяснение состоит в том, что у случайных данных спектр
ковариационной матрицы входов близок к вырожденному по масштабу (все
собственные значения одного порядка), тогда как у digits после проекции на
главные компоненты спектр существенно неоднороден. Многомасштабность целевой
функции --- именно тот фактор, который по результатам~\cite{kong2020}
порождает хаотическое поведение детерминированного градиентного спуска.
Фрактальность границы, таким образом, не является артефактом синтетических
данных.

\subsection{Трёхмерное пространство $(\eta,\beta,\lambda_{wd})$}\label{sec:3d}

Трёхмерный эксперимент реализован как набор двумерных сечений
$(\eta,\beta)$ при четырёх значениях затухания весов
(табл.~\ref{tab:wd}, рис.~\ref{fig:wdslices}).

\begin{table}[htbp]
  \centering
  \caption{Сечения трёхмерного пространства $(\eta,\beta,\lambda_{wd})$.}
  \label{tab:wd}
  \small
  \input{tables/tab_wd_slices.tex}
\end{table}

\begin{figure}[htbp]
  \centering
  \includegraphics[width=\linewidth]{f13_wd_slices.png}
  \caption{Сечения $(\eta,\beta)$ при разных значениях $\lambda_{wd}$.}
  \label{fig:wdslices}
\end{figure}

Зависимость немонотонна: $\lambda_{wd}=0$ даёт $\DWdZero$, $10^{-4}$ ---
$\DWdOne$, $10^{-3}$ --- $\DWdTwo$, но $10^{-2}$ --- снова $\DWdThree$.
Содержательно это означает, что умеренная регуляризация сглаживает границу,
а при дальнейшем росте $\lambda_{wd}$ появляется новый эффект: затухание весов
начинает конкурировать с обучением, и возникает дополнительная граница между
режимом обучения и режимом стягивания весов к нулю, которая сама по себе
нерегулярна. Заметно также, что доля осцилляционного режима резко возрастает
(с $0.20$ до $0.48$).

Отсюда следует прямой ответ на вопрос, поставленный в задании: по набору
двумерных сечений судить о сложности границы в полном трёхмерном пространстве
можно лишь с осторожностью. Немонотонность $D(\lambda_{wd})$ показывает, что
сечения не являются «типичными» срезами однородного объекта; каждое сечение
характеризует локальную геометрию границы вблизи выбранной плоскости, а не
глобальную структуру. Строгое утверждение о размерности трёхмерной границы
потребовало бы прямого box-counting в трёхмерной сетке, что превышает
доступные вычислительные ресурсы: сетка $256^3$ при $T=400$ потребовала бы
примерно в 256 раз больше времени, чем одно двумерное сканирование.

\subsection{Последовательность зумов}

Пять уровней зума вокруг автоматически выбранной точки границы
($\eta^\ast\approx0.369$, $\beta^\ast\approx0.300$) показаны на
рис.~\ref{fig:zoom}, измеренные размерности --- в табл.~\ref{tab:zoom}.

\begin{table}[htbp]
  \centering
  \caption{Размерность границы на разных уровнях зума в плоскости
  $(\eta,\beta)$.}
  \label{tab:zoom}
  \small
  \input{tables/tab_zoom.tex}
\end{table}

\begin{figure}[htbp]
  \centering
  \includegraphics[width=\linewidth]{f14_zoom.png}
  \caption{Последовательность зумов: каждый следующий уровень сужает окно
  вдвое по $\log\eta$ и по $\beta$.}
  \label{fig:zoom}
\end{figure}

Структура сохраняется на всех пяти уровнях: качественный характер картины не
меняется, что и означает самоподобие. Однако, в отличие от предыдущей НИР, где
размерность по уровням зума была практически постоянной
($\sigma = 0.009$), здесь наблюдается систематический рост от
$1.679$ на нулевом уровне до $1.855$ на четвёртом с последующим насыщением;
медиана составляет $\DZoomMedian$ при стандартном отклонении $\DZoomStd$.

Причина роста видна из долей режимов: с углублением зума область переполнения
исчезает полностью, а доля осцилляционного режима возрастает с $0.26$ до
$0.63$. Иными словами, при увеличении мы попадаем внутрь зоны, где режимы
сходимости и колебаний перемешаны на масштабе отдельных узлов сетки, и
измеряемое множество перестаёт быть кривой, приближаясь к площадному
заполнению ($D\to2$). Это не столько свойство границы, сколько указание на
предел применимости самой процедуры: box-counting по бинаризованной карте
измеряет геометрию перемешанной зоны, а не гладкой кривой. Рис.~\ref{fig:selfsim}
показывает извлечённые границы на нулевом и четвёртом уровнях.

\begin{figure}[htbp]
  \centering
  \includegraphics[width=\linewidth]{f16_boundary_selfsim.png}
  \caption{Извлечённая граница на нулевом и четвёртом уровнях зума.}
  \label{fig:selfsim}
\end{figure}

\subsection{Влияние численной точности}

Сравнение расчёта в \texttt{float64} и \texttt{float32} для одного и того же
сечения $(\eta,\beta)$ приведено в табл.~\ref{tab:prec} и на
рис.~\ref{fig:prec}.

\begin{table}[htbp]
  \centering
  \caption{Влияние типа точности на оценку размерности.}
  \label{tab:prec}
  \small
  \begin{tabular}{lccccc}
\toprule
Конфигурация & $D_{\mathrm{box}}$ & сход. & осц. & расх. & NaN\\
\midrule
float64 & $1.616 \pm 0.042$ & 0.52 & 0.20 & 0.05 & 0.23\\
float32 & $1.716 \pm 0.060$ & 0.44 & 0.28 & 0.00 & 0.28\\
\bottomrule
\end{tabular}

\end{table}

\begin{figure}[htbp]
  \centering
  \includegraphics[width=\linewidth]{f15_precision.png}
  \caption{Одно и то же сечение $(\eta,\beta)$, вычисленное в двойной и
  одинарной точности.}
  \label{fig:prec}
\end{figure}

Крупномасштабная структура совпадает, что подтверждает: наблюдаемая
изрезанность не является артефактом численной ошибки. Оценка размерности в
одинарной точности оказалась выше ($\DFloatThirtyTwo$ против $\DWdZero$).
Это ожидаемо: в хаотической зоне траектории чувствительны к начальным
условиям, и ошибки округления в \texttt{float32} действуют как дополнительный
шум, увеличивая перемешивание режимов вблизи границы. Для измерения
размерности двойная точность необходима.
